% Technical requirements slides (v0.9)

\begin{frame}{Рабочие места участников\orient}
  \vspace{0.1em}
  {\footnotesize\color{acGray}Требования к рабочим станциям на очных площадках
  всех этапов \textbf{совпадают}.}
  \vspace{0.6em}

  \small
  \begin{tabular}{@{}p{0.19\textwidth}p{0.32\textwidth}p{0.43\textwidth}@{}}
    \toprule
    \textbf{Критерий} & \textbf{Минимальные} & \textbf{Рекомендуемые}\\
    \midrule
    Процессор & Intel Core i5 / AMD Ryzen 5 (или аналог) &
      Intel Core i5, i7 12-го поколения или аналог: 6+ ядер, 12+ потоков, 3{,}0/4{,}6 ГГц\\
    \addlinespace[0.25em]
    ОЗУ & 16 ГБ DDR4 / DDR5 (двухканальная конфигурация) &
      16 ГБ и более, DDR5 4800+ МГц, двухканальная конфигурация\\
    \addlinespace[0.25em]
    Накопитель & SSD 250 ГБ &
      SSD 500+ ГБ, M.2, 3D TLC NAND, 3100/2600 МБ/с и выше, 400\,000/470\,000 IOPS и выше\\
    \addlinespace[0.25em]
    Сеть & Ethernet 100 Мбит/с & Ethernet 1 Гбит/с\\
    \addlinespace[0.25em]
    ОС & Linux x64 (Ubuntu 26.04 LTS или аналог), Windows 10/11 &
      Linux x64 (Ubuntu 26.04 LTS или аналог), Windows 10/11\\
    \bottomrule
  \end{tabular}
\end{frame}

\begin{frame}{Оборудование для учебно-тренировочных сборов\orient}
  \vspace{0.1em}
  {\footnotesize\color{acGray}Отличие от рабочих мест этапов ---
  графический ускоритель. Потребность: \textbf{25 устройств}.}
  \vspace{0.6em}

  \small
  \begin{tabular}{@{}p{0.22\textwidth}p{0.34\textwidth}p{0.38\textwidth}@{}}
    \toprule
    \textbf{Критерий} & \textbf{Минимальные} & \textbf{Рекомендуемые}\\
    \midrule
    Процессор & Intel Core i5 / AMD Ryzen 5 (или аналог) &
      Intel Core i5, i7 12-го поколения или аналог: 6+ ядер, 12+ потоков, 3{,}0/4{,}6 ГГц\\
    \addlinespace[0.25em]
    ОЗУ & 16 ГБ DDR4 / DDR5 (двухканальная конфигурация) &
      16 ГБ и более, DDR5 4800+ МГц\\
    \addlinespace[0.25em]
    \textbf{Видеокарта} & \textbf{NVIDIA GeForce RTX 5070 8 ГБ} (или аналог) &
      \textbf{NVIDIA GeForce RTX 5090 24 ГБ} (или аналог)\\
    \addlinespace[0.25em]
    Накопитель & SSD 250 ГБ & SSD 500+ ГБ, M.2, 3D TLC NAND\\
    \bottomrule
  \end{tabular}
\end{frame}

\begin{frame}{Среда выполнения: обязательные пакеты\orient}
  \vspace{0.2em}
  {\footnotesize\color{acGray}Пакеты должны быть установлены в окружение
  рабочей машины и присутствовать в тестирующей среде.}
  \vspace{0.6em}

  \footnotesize
  \begin{tabular}{@{}p{0.22\textwidth}p{0.74\textwidth}@{}}
    \toprule
    \textbf{Категория} & \textbf{Пакеты}\\
    \midrule
    Основные библиотеки &
      \ttfamily torch, torchvision, torchaudio, transformers, accelerate, peft, trl,
      scikit-learn, xgboost, lightgbm, catboost, sentence-transformers, datasets,
      evaluate, spacy, nltk, gensim, fasttext\\
    \addlinespace[0.3em]
    Обработка данных &
      \ttfamily numpy, pandas, scipy, polars, pyarrow, h5py\\
    \addlinespace[0.3em]
    Компьютерное зрение &
      \ttfamily opencv-python, Pillow, scikit-image, albumentations\\
    \addlinespace[0.3em]
    Визуализация данных &
      \ttfamily matplotlib, seaborn, plotly\\
    \addlinespace[0.3em]
    Утилиты и обучение &
      \ttfamily tqdm, joblib, tensorboard, pytorch-lightning, pydantic, pyyaml\\
    \bottomrule
  \end{tabular}
\end{frame}

\begin{frame}{Сеть на площадках III и IV этапов\orient}
  \vspace{0.3em}
  \begin{columns}[T]
    \begin{column}{0.49\textwidth}
      \begin{block}{Обязательные условия}
        \begin{itemize}\setlength\itemsep{0.3em}
          \item Оргкомитет этапа имеет \textbf{непосредственный доступ
                к настройкам маршрутизатора и сетевого оборудования} ---
                для самостоятельного введения ограничений.
          \item Общая пропускная способность внешнего канала ---
                \textbf{не менее 200 Мбит/с}.
        \end{itemize}
      \end{block}
    \end{column}
    \begin{column}{0.49\textwidth}
      \begin{block}{Что это означает на практике}
        \begin{itemize}\setlength\itemsep{0.3em}
          \item площадка выделяется в управляемый сегмент, не зависящий
                от общей сети учреждения;
          \item доступ в Интернет ограничивается за пределами АТС
                (белый список ресурсов);
          \item права администрирования на время туров передаются
                техническому представителю оргкомитета;
          \item до тура проводится тестовый прогон площадки.
        \end{itemize}
      \end{block}
    \end{column}
  \end{columns}
\end{frame}
